% docs/texte/90_appendix/fringe.tex -- appendix chapter "The Term Fringe" (Interferosim)
% Version: 0.1.0 (2026-10-12)
% Detailed explanation of the term "fringe" and the related terms used in
% the code and documentation (fringe spacing, rate, contrast, stopping,
% fitting, washing). Notation as in texte/05_theory/notation.tex.
% ================================================================

\index{fringe}
\section{Origin of the term}

\emph{Fringe} (German \emph{Interferenzstreifen}) comes from optics. When
two coherent light beams overlap, the screen shows alternating bright
and dark bands -- the interference fringes of Young's double-slit
experiment. Michelson and Pease used the disappearance of such fringes
to measure the first stellar diameter, that of Betelgeuse
\cite{michelson1921}. Radio astronomy took over the term together with
the method: Ryle and Vonberg built the first radio interferometer to
observe the Sun \cite{ryle1946}. In a radio interferometer there is no
screen; the fringes appear in the correlator output
(Fig.~\ref{fig:two-element}).

\section{Three views of the same fringe}

The response of a two-element interferometer to a point source,
Eq.~\eqref{eq:fringe}, can be read in three equivalent ways.

\paragraph*{As a pattern on the sky.} With $\omega\tau_g =
2\pi\,(b/\lambda)\cos\theta$ the normalised response is
\begin{equation}
  \frac{R(\theta)}{R_{\max}} = \cos\!\left(2\pi\,\frac{b}{\lambda}\cos\theta\right).
\end{equation}
The phase changes by $2\pi$ when $\cos\theta$ changes by $\lambda/b$, so
the angular fringe spacing is
\begin{equation}
  \Delta\theta = \frac{\lambda}{b\,\sin\theta}
  \;\xrightarrow{\ \theta=90^\circ\ }\; \frac{\lambda}{b} .
\end{equation}
The interferometer behaves like a grating of alternating positive and
negative sensitivity laid over the sky (Fig.~\ref{fig:fringe-pattern}).
Towards the direction of the baseline ($\theta\to0^\circ, 180^\circ$)
the fringes become wider, because the projected baseline
$b\sin\theta$ becomes shorter.

\paragraph*{As a signal in time.} The Earth's rotation moves the source
through this pattern. The output oscillates with the fringe frequency
$f_F$ of Eq.~\eqref{eq:fringerate}. This is how the fringes are seen in
a measurement: as a slow oscillation of the correlator output (the
\emph{fringe rate}), not as a picture.

\paragraph*{As a Fourier component.} A baseline $(u,v)$ responds only to
the component $e^{-2\pi i(ul+vm)}$ of the sky brightness
(Eq.~\eqref{eq:vis}). The fringe pattern \emph{is} this Fourier
component. Amplitude and phase of the fringe are amplitude and phase
of the complex visibility $\mathcal V(u,v)$; aperture synthesis
collects many such fringes and adds them up to an image.

\begin{figure}[htbp]
  \centering
  % docs/texte/90_appendix/figures/fringe_pattern.tex -- fringe pattern of a two-element interferometer (pgfplots)
% Version: 0.1.1 (2026-10-12) -- legend above the plot
% Normalised response R(theta) = cos(2 pi (b/lambda) cos theta) for b = 10 lambda,
% monochromatic and with a fractional bandwidth of 10 % (envelope
% |sinc(dnu tau_g)|, dnu tau_g = (dnu/nu)(b/lambda) cos theta).
\begin{tikzpicture}
  \begin{axis}[
      width=0.95\textwidth, height=6cm,
      xmin=0, xmax=180, ymin=-1.15, ymax=1.15,
      xtick={0,30,60,90,120,150,180},
      xlabel={angle $\theta$ between baseline and source direction [deg]},
      ylabel={normalised response $R/R_{\max}$},
      legend style={font=\footnotesize, at={(0.5,1.12)}, anchor=south, draw=none},
      legend columns=3,
      tick label style={font=\footnotesize}, label style={font=\footnotesize},
      samples=1200, domain=0:180]
    % monochromatic
    \addplot[thin, gray] {cos(deg(2*pi*10*cos(x)))};
    \addlegendentry{monochromatic}
    % 10 % bandwidth: fringe times sinc envelope
    \addplot[thick, blue!70!black]
      {cos(deg(2*pi*10*cos(x))) * sin(deg(pi*0.1*10*cos(x)+1e-9))/(pi*0.1*10*cos(x)+1e-9)};
    \addlegendentry{$\Delta\nu/\nu = 0.1$}
    \addplot[thin, dashed, red!70!black]
      {abs(sin(deg(pi*0.1*10*cos(x)+1e-9))/(pi*0.1*10*cos(x)+1e-9))};
    \addlegendentry{envelope $|\operatorname{sinc}(\Delta\nu\,\tau_g)|$}
    % mark fringe spacing near theta = 90 deg
    \draw[<->, thick] (axis cs:90,1.07) -- (axis cs:95.73,1.07)
      node[midway, above, font=\scriptsize] {$\approx\lambda/b$};
  \end{axis}
\end{tikzpicture}

  \caption{Fringe pattern of a baseline $b = 10\,\lambda$ as a function
           of the source direction. Grey: monochromatic; blue: with a
           fractional bandwidth of 10\,\%, the fringes are confined to
           an envelope around $\theta = 90^\circ$ (bandwidth smearing,
           \enquote{delay beam}).}
  \label{fig:fringe-pattern}
\end{figure}

\section{Fringe visibility and source size}

Michelson defined the visibility of the fringes as the contrast
\begin{equation}
  \frac{I_{\max}-I_{\min}}{I_{\max}+I_{\min}} ,
\end{equation}
with the maximum and minimum intensity of the fringe pattern. For an
adding interferometer this is the normalised visibility
$|\mathcal V|/S$; the name \enquote{visibility} for $\mathcal V$ has its
origin here. A point source gives full contrast. An extended source
is a superposition of point sources whose fringes are shifted against
each other, so the contrast decreases. For a uniform disc it vanishes
at $b_{\mathrm{null}} = 1.2197\,\lambda/\theta_d$ -- this is the
principle of Michelson's diameter measurement and of the solar
diameter measurement in notebook~B.

\section{Related terms}

\begin{description}
  \item[Fringe spacing] Angular period $\lambda/b$ of the pattern; the
        resolution scale of a baseline.
  \item[Fringe rate] Temporal frequency $f_F$ of the fringes caused by
        the Earth's rotation, Eq.~\eqref{eq:fringerate}.
  \item[Fringe amplitude, fringe phase] Amplitude and phase of the
        correlator output; together the complex visibility.
  \item[Fringe stopping] Removing the predicted fringe rate by a phase
        rotation of one signal, so that the correlator output varies
        only slowly and can be averaged over long intervals.
  \item[Delay tracking] Compensating the geometric delay $\tau_g$, so
        that the fringes stay in the centre of the bandwidth envelope.
  \item[Fringe fitting / fringe search] Determining the residual delay
        and rate when they are not known in advance (independent clocks
        in VLBI); the correlation is searched over a grid of delays and
        rates.
  \item[Fringe washing (bandwidth smearing)] Loss of fringe amplitude
        by the factor $|\operatorname{sinc}(\Delta\nu\,\tau)|$ when
        different frequencies in the band produce different fringe
        phases (Fig.~\ref{fig:fringe-pattern}); in time the same effect
        occurs if the averaging interval is not short compared with
        $1/f_F$ (time smearing).
\end{description}

\section{Numbers for the Interferosim examples}

\begin{center}
\begin{tabular}{@{}lccc@{}}
\toprule
Baseline & $b/\lambda$ & fringe spacing $\lambda/b$ & fringe period at transit ($\delta=0$) \\
\midrule
10\,m (horn test, Sun) & 47 & $1.21^\circ$ & 290\,s \\
11.5\,km (long baseline) & 54\,500 & $3.8''$ & -- \\
\bottomrule
\end{tabular}
\end{center}
At $\lambda = 21.1\,\mathrm{cm}$. The fringe period of the 10\,m
baseline uses the solar rotation rate (notebook~B).
